% ---------------------------------------------------------------------
\section{Introduction}\label{sec:intro}
% ---------------------------------------------------------------------

A charge point operator (CPO) buys energy for its chargers at time-varying prices, and vehicle-to-grid (V2G) discharging lets the same operator sell energy from parked electric vehicles (EVs) when prices are high.
Two hard requirements limit this freedom.
The aggregate power of the station must stay within the transformer limit, and each EV must hold its requested energy when it leaves.
A schedule that breaks the first requirement overloads equipment, and one that breaks the second requirement breaks the service promise to the driver.

Reinforcement learning (RL) has become a common approach to this control task, because prices, arrivals and departures are uncertain and hard to model in advance \cite{qiu2023reinforcement,michailidis2025reinforcement}.
Most RL controllers for EV charging encode the departure energy requirement and the transformer limit as weighted penalties in the reward \cite{sadeghianpourhamami2020definition,ye2022learning,korkas2024distributed,orfanoudakis2025ev2gym}.
A penalty trades violations against profit, so a trained policy may still violate either requirement at test time \cite{orfanoudakis2025ev2gym,tabas2022computationally}.
Constrained formulations replace the fixed weight with a bound on the expected cumulative cost \cite{altman2021constrained}, but such a bound holds on average over episodes and not at every step \cite{gu2024review,su2025review}.

A second limitation concerns station size.
The default interface of the EV2Gym simulator assigns one action to each charger and uses a flat observation, so a trained network is tied to one number of chargers \cite{orfanoudakis2025ev2gym}.
Graph and attention policies remove this restriction \cite{orfanoudakis2025scalable,orfanoudakis2026graph,wang2026hetgps}, yet in these size-agnostic controllers the constraints still enter through reward terms, imitation data or heuristic filters.
The authors of one of them state explicitly that their filter makes no pointwise feasibility claim \cite{wang2026hetgps}.

Safety layers correct each action before execution and can enforce limits by construction \cite{dalal2018safe,chen2021enforcing,tabas2022computationally}.
For EV stations, existing layers cover per-step limits such as transformer capacity and charger current \cite{wan2025logicinformed}, or they apply rule-based shields \cite{guan2024rule,zhang2025dumes}.
A per-step layer alone does not protect the departure energy requirement, because that requirement couples the present action with every later step of the stay.
Per-EV latest-start bounds address this coupling in laxity-based scheduling \cite{chen2022smoothed} and in hierarchical V2G scheduling \cite{yan2026hierarchical}.
To our knowledge, no published controller combines such bounds with a learned policy that runs at any station size and comes with a proof of multi-step feasibility.
\TODO{Confirm this sentence against the full texts of \cite{orfanoudakis2026can}, \cite{song2025longterm} and \cite{eland2026safe} before submission (novelty risks listed in the gap analysis)}.

This paper proposes a safe attention policy (\method) for V2G charging control at a CPO station.
An attention policy treats each charger as a token and outputs one raw power action per token with shared weights.
A safety layer, a closed-form projection onto the feasible set, maps the raw actions to commanded powers that respect the per-EV bounds and the transformer limit.
The projection has the same form as the rate rule of smoothed least-laxity-first (sLLF) scheduling \cite{chen2022smoothed}, and it is the classical Euclidean projection onto a box intersected with one sum constraint \cite{kiwiel2007breakpoint,wang2022efficient}.
Our contribution therefore lies in the bounds fed to this projection, in its use during training and deployment, and in the multi-step argument built on it.

The paper makes four contributions, and each of them corresponds to a part of the evaluation in \Sec{sec:results}.
\begin{enumerate}
\item We design an attention policy for V2G charging control in which each charger is a token and a shared head produces one action per token. The policy is permutation-equivariant, ignores empty chargers by masking and has a parameter count that does not depend on the number of chargers.
\item We derive a safety layer whose per-EV lower bound follows from a reachable-energy operator that includes the charging taper. Theorem~\ref{thm:invariance} proves that the transformer limit and every departure energy requirement hold at every step for any raw action, under stated assumptions. A least-laxity-first fallback rule and admission control handle the steps on which aggregate feasibility fails.
\item We define a zero-shot evaluation in EV2Gym that trains on 10 and 20 chargers and tests on up to 200 chargers without retraining. Performance is measured as the optimality gap to a perfect-information linear program with hard constraints, and the cost of the guarantee is measured against penalty-based training of the same network.
\item As a secondary contribution, we replace the known-departure assumption with a split-conformal lower bound on the departure step, and Proposition~\ref{prop:conformal} turns the guarantee into a coverage-level statement.
\end{enumerate}
\Sec{sec:related} shows that none of these ingredients is new in isolation, so the claim of the paper rests on their combination and on the measured price of the guarantee.

\Sec{sec:related} reviews related work, and \Sec{sec:problem} states the system model and the Markov decision process (MDP).
The attention policy, the safety layer and the training algorithm follow in \Sec{sec:method}, and \Sec{sec:setup} describes the experimental setup.
Results appear in \Sec{sec:results} and the discussion in \Sec{sec:discussion}.
\Sec{sec:conclusion} concludes the paper, and the supplementary material contains all proofs, the EV2Gym modelling details, the training details and further result tables.
